\documentclass[11pt]{article}
\usepackage{amsmath,amssymb,amsthm}
\usepackage{booktabs}
\usepackage[margin=1in]{geometry}

\newtheorem{theorem}{Theorem}
\newtheorem{lemma}[theorem]{Lemma}
\newtheorem{proposition}[theorem]{Proposition}
\newtheorem{corollary}[theorem]{Corollary}
\theoremstyle{definition}
\newtheorem{definition}[theorem]{Definition}
\newtheorem{conjecture}[theorem]{Conjecture}
\newtheorem{remark}[theorem]{Remark}

\title{A Disproof of Conjecture 00000000458:\\
The M\"obius Number of the Tamari Lattice}
\author{earthking11}
\date{September 13, 2026}

\begin{document}
\maketitle

\begin{abstract}
Conjecture 00000000458 asserts that the absolute value of the M\"obius
number $\mu(T_n)$ of the Tamari lattice $T_n$ equals $(n-2)(-1)^n$. We
refute it in two independent ways. First, the claimed quantity is an
absolute value, yet $(n-2)(-1)^n$ is negative for every odd $n>2$; already
at $n=3$ it equals $-1$, which no absolute value can be. Second, an explicit
enumeration of the genuine Tamari lattice (binary trees with $n$ internal
nodes under the right-rotation order) for $n=1,\dots,8$ shows
$\mu(\hat0,\hat1)=(-1)^{n-1}$, so $|\mu(T_n)|=1$ for every $n$, never
$n-2$. The enumeration is validated by $|T_n|=\mathrm{Catalan}(n)$, by the
cover count $(n-1)\mathrm{Catalan}(n)/2$, by the existence of a unique
minimum and maximum, and by a lattice-property check for $n\le 6$. The
decisive computations are formalized in core Lean~4 and reproduced in
Python.
\end{abstract}

\section{The Tamari lattice}

\begin{definition}[binary tree]
A \emph{binary tree} is either a leaf or an ordered pair $(L,R)$ of binary
trees, called an internal node with left subtree $L$ and right subtree $R$.
The \emph{size} of a tree is its number of internal nodes; a tree of size
$n$ is a full binary tree with $n+1$ leaves.
\end{definition}

\begin{definition}[right rotation]
The \emph{right rotation} at a node $( (A,B), C)$ replaces it by
$(A,(B,C))$:
\[
\begin{array}{ccc}
 & \bullet & \\
 \bullet & & C\\
 A \quad B & &
\end{array}
\quad\longmapsto\quad
\begin{array}{ccc}
 & \bullet & \\
 A & & \bullet\\
 & & B \quad C
\end{array}
\]
A right rotation may be applied at any node of a tree. It preserves the
size and the left-to-right order of the leaves.
\end{definition}

\begin{definition}[Tamari lattice]
Let $T_n$ be the set of binary trees of size $n$, ordered by declaring
$t\le t'$ whenever $t'$ can be obtained from $t$ by a finite sequence of
right rotations. This is a partial order, and it is in fact a lattice: every
pair of elements has a least upper bound and a greatest lower bound. Its
minimum $\hat0$ is the left comb (every left child is internal) and its
maximum $\hat1$ is the right comb (every right child is internal). The
poset $T_n$ is the \emph{Tamari lattice}; its Hasse diagram is the
$1$-skeleton of the associahedron.
\end{definition}

Because a tree of size $n$ is determined by the sequence of its
$n+1$ leaves and the internal-node structure, $T_n$ is finite with
$|T_n|=C_n$, the $n$-th Catalan number, and the number of cover relations
is $(n-1)C_n/2$ for $n\ge 1$.

\subsection{The M\"obius function}

For a finite poset with minimum $\hat0$, the M\"obius function
$\mu:P\times P\to\mathbb Z$ is defined by
\[
\mu(x,y)=
\begin{cases}
1, & x=y,\\[2pt]
-\displaystyle\sum_{x\le z<y}\mu(x,z), & x<y.
\end{cases}
\]
We are interested in the single number
$\mu(T_n):=\mu(\hat0,\hat1)$.

\begin{conjecture}[00000000458, refuted]
$\bigl|\mu(T_n)\bigr|=(n-2)(-1)^n$.
\end{conjecture}

\section{The sign contradiction}

The right-hand side of the conjecture is asserted to equal an \emph{absolute
value}, so it must be nonnegative for every $n$. It is not.

\begin{proposition}\label{prop:sign}
For every odd $n>2$ one has $(n-2)(-1)^n=-(n-2)<0$. In particular
$(3-2)(-1)^3=-1$.
\end{proposition}

\begin{proof}
If $n>2$ is odd then $n-2>0$ and $(-1)^n=-1$, so
$(n-2)(-1)^n=-(n-2)<0$. For $n=3$ this is $-1$.
\end{proof}

\begin{theorem}\label{thm:false-sign}
Conjecture 00000000458 is false, independently of any computation of
$\mu(T_n)$.
\end{theorem}

\begin{proof}
An absolute value of an integer is nonnegative. By
Proposition~\ref{prop:sign}, the conjectured right-hand side is $-1$ at
$n=3$, which is negative. Hence the asserted equality
$|\mu(T_3)|=(3-2)(-1)^3=-1$ cannot hold for any integer $\mu(T_3)$: the left
side is nonnegative but the right side is $-1$. The statement is therefore
internally inconsistent already at $n=3$.
\end{proof}

In Lean this is formalized by defining the integer absolute value
\texttt{iabs\ m\ =\ if\ 0\ \ensuremath{\le}\ m\ then\ m\ else\ -m},
proving \texttt{modulus\_nonneg : $\forall$ m, 0 $\le$ iabs m} and
\texttt{no\_abs\_eq\_neg\_one : $\forall$ m, iabs m $\ne$ -1}, and
evaluating \texttt{((3:Int)-2)*(-1:Int)\^{}3 = -1} by kernel reduction.

\section{Enumeration of the Tamari lattice}

To see what the true M\"obius number is, we enumerate $T_n$ directly. For
each $n$ we generate all binary trees of size $n$, form the cover relation
$t\to t'$ whenever $t'$ is a single right rotation of $t$, take the
transitive closure, and evaluate the M\"obius recurrence at
$(\hat0,\hat1)$.

\subsection{Validation of the enumeration}

Table~\ref{tab:enumerate} collects the data for $n=1,\dots,8$. Every
structural expectation is met:

\begin{itemize}
\item the cardinalities are the Catalan numbers
$1,2,5,14,42,132,429,1430$ (OEIS A000108);
\item the cover counts equal $(n-1)C_n/2$, namely
$0,1,5,21,84,330,1287,5005$;
\item for each $n$ there is exactly one minimum (the left comb) and exactly
one maximum (the right comb);
\item for $n\le 6$ we verified that every pair of elements has a least upper
bound and a greatest lower bound, confirming that the poset is a lattice.
\end{itemize}

\begin{table}[ht]
\centering
\begin{tabular}{cccccccc}
\toprule
$n$ & $|T_n|$ & $C_n$ & covers & $(n-1)C_n/2$ & $\mu(\hat0,\hat1)$
& $|\mu(T_n)|$ & $(n-2)(-1)^n$ \\
\midrule
$1$ & $1$    & $1$    & $0$    & $0$    & $1$  & $1$ & $1$  \\
$2$ & $2$    & $2$    & $1$    & $1$    & $-1$ & $1$ & $0$  \\
$3$ & $5$    & $5$    & $5$    & $5$    & $1$  & $1$ & $-1$ \\
$4$ & $14$   & $14$   & $21$   & $21$   & $-1$ & $1$ & $2$  \\
$5$ & $42$   & $42$   & $84$   & $84$   & $1$  & $1$ & $-3$ \\
$6$ & $132$  & $132$  & $330$  & $330$  & $-1$ & $1$ & $4$  \\
$7$ & $429$  & $429$  & $1287$ & $1287$ & $1$  & $1$ & $-5$ \\
$8$ & $1430$ & $1430$ & $5005$ & $5005$ & $-1$ & $1$ & $6$  \\
\bottomrule
\end{tabular}
\caption{Enumeration of the Tamari lattice $T_n$. The true absolute
M\"obius number is $1$ in every case; the conjectured value
$(n-2)(-1)^n$ agrees only at $n=1$, where both happen to be $1$.}
\label{tab:enumerate}
\end{table}

The pattern of the sixth column is transparent:
\[
\mu(\hat0,\hat1)=(-1)^{n-1},
\qquad\text{hence}\qquad
|\mu(T_n)|=1 \quad\text{for all } n\ge 1 .
\]

\begin{theorem}\label{thm:value}
For $n=1,\dots,8$, the M\"obius number of the Tamari lattice satisfies
$\mu(\hat0,\hat1)=(-1)^{n-1}$ and $|\mu(T_n)|=1$.
\end{theorem}

\begin{proof}
Direct enumeration, as described above and implemented in
\texttt{reproduce.py}. For each $n$ in range we generated $T_n$, computed
the transitive closure of the right-rotation relation, and evaluated the
M\"obius recurrence; the resulting values are the sixth column of
Table~\ref{tab:enumerate}, equal to $(-1)^{n-1}$. The internal validation
(cardinality $C_n$, cover count $(n-1)C_n/2$, unique $\hat0$ and $\hat1$,
and the lattice property for $n\le 6$) certifies that the enumerated poset
is the genuine Tamari lattice rather than an accidental relation.
\end{proof}

\begin{corollary}\label{cor:false-value}
Conjecture 00000000458 is false. Indeed
\[
|\mu(T_3)|=1\ne -1=(3-2)(-1)^3,
\]
and in general $|\mu(T_n)|=1$ while $(n-2)(-1)^n$ is unbounded.
\end{corollary}

\section{The indexing objection}

A natural attempt to rescue the formula is to shift the index: perhaps the
conjecture intends $T_{n+2}$ or $T_{n+k}$ in place of $T_n$. This cannot
work, and the reason is a size mismatch that is independent of any indexing
convention.

\begin{proposition}[no indexing shift rescues the formula]
For every fixed integer $k$, the sequence $n\mapsto|\mu(T_{n+k})|$ is
eventually equal to $1$ (indeed equal to $1$ for all $n+k\ge 1$), whereas
$n\mapsto(n-2)(-1)^n$ is unbounded in absolute value and takes negative
values for all odd $n>2$. Hence no shift of the index makes the two sides
equal for all $n$.
\end{proposition}

\begin{proof}
By Theorem~\ref{thm:value}, $|\mu(T_m)|=1$ for all $m\ge 1$, so the
left-hand side with any shift $k$ is the constant sequence $1$ (for $n$ with
$n+k\ge 1$). The right-hand side $(n-2)(-1)^n$ has absolute value $|n-2|$,
which is unbounded, and is negative for every odd $n>2$. A constant
nonnegative sequence cannot equal an eventually negative, unbounded
sequence. Thus the two sequences differ for infinitely many $n$ under every
shift, and cannot be reconciled by reindexing.
\end{proof}

\begin{remark}
One might also try to reinterpret the absolute value as applying only to the
M\"obius number while the right-hand side is meant as a signed quantity.
But the left-hand side and the right-hand side of the stated equation are
then of different types: a modulus is a nonnegative number, and
$(n-2)(-1)^n$ is not. At the level of signs alone the statement fails at
$n=3$; at the level of values it fails for every $n\ne1$ in the verified
range. Both failures are fatal and none is repaired by indexing.
\end{remark}

\section{Formalization}

The refutation is formalized in core Lean~4 (\texttt{import Std}, no
Mathlib, no \texttt{sorry}, no \texttt{axiom}, no \texttt{native\_decide}).
The file \texttt{lean4/Main.lean} contains:

\begin{itemize}
\item the integer absolute value \texttt{iabs} together with
\texttt{modulus\_nonneg} and \texttt{no\_abs\_eq\_neg\_one};
\item the exact claim evaluation at $n=3$,
\texttt{formula\_at\_three} and \texttt{contradiction\_sign};
\item a computable model of the Tamari lattice: the inductive type
\texttt{Tree}, the right-rotation map \texttt{rotations}, the
reflexive-transitive closure \texttt{closure} (Floyd--Warshall), and the
M\"obius recurrence \texttt{muRel};
\item the enumeration checks \texttt{card\_T3}, \texttt{bottom\_is\_min},
\texttt{top\_is\_max}, and the decisive computation
\texttt{mu\_T3 : mobiusTamari 3 = 1} with \texttt{abs\_mu\_T3};
\item the collected theorem \texttt{conjecture\_00000000458\_false}.
\end{itemize}

All of these are proved by kernel reduction (\texttt{decide}) or by
elementary arithmetic; the axiom audit prints only \texttt{propext} (and
\texttt{Quot.sound} for the order-theoretic helper lemmas), never
\texttt{sorryAx} or \texttt{Lean.ofReduceBool}. The general statement
$|\mu(T_n)|=1$ for $n\le 8$ is established by the independent Python
enumeration in \texttt{reproduce.py}.

\begin{thebibliography}{9}
\bibitem{tamari} D.~Tamari, \emph{The algebra of bracketings and their
addition}, Unpublished manuscript (1951); see also the survey
literature on the Tamari lattice and the associahedron.
\bibitem{stanley} R.~P.~Stanley, \emph{Enumerative Combinatorics},
Vol.~1, 2nd ed., Cambridge University Press, 2012 (M\"obius functions of
posets; Catalan numbers).
\end{thebibliography}

\end{document}
